\section{Meta-models: why a fast model is built from a slow one}
A detailed model such as Daisy takes minutes or hours for a run of twenty years, which is too slow for a decision tool in which the user moves a slider and expects an answer at once. A \kw{meta-model} (surrogate, emulator) is a simple model that is fitted to the results of the detailed one, and then replaces it inside the tool. Two ways are used. The first is a \emph{statistical emulator}, a regression or a neural network of the outputs on the inputs; it is fast but gives no insight and fails outside the range of the training data. The second, used here, is a \emph{physically structured meta-model}: a water balance with few parameters in which the results of Daisy supply the ratios and the scaling factors. It keeps the conservation of water by construction, and the parameters have a meaning.
\par\noindent \textbf{Note:} A meta-model must be tested on runs that were not used to fit it (held-out runs), and it must say when the user is outside the range for which it was fitted.

\section{The bucket model of the soil water}
The simplest model of the soil water treats the root zone as a \kw{bucket}: it is filled by rain and irrigation, emptied by evapotranspiration, and overflows as drainage when it is full. In daily steps the storage $S$ is
\begin{equation*}\boxed{S_{d+1}=\min\bigl(S_{max},\ S_d+P_d+I_d-ET_{a,d}\bigr),\qquad D_d=\max\bigl(0,\ S_d+P_d+I_d-ET_{a,d}-S_{max}\bigr)}\end{equation*}
where $S_{max}$ is the capacity of the bucket (field capacity), $P$ the rain, $I$ the irrigation, $ET_a$ the actual evapotranspiration and $D$ the drainage. With a second bucket for the layer below it is called a \kw{tipping-bucket} (cascade) model. The model has two advantages (it conserves water exactly and needs few parameters) and one limit: it has no suction, so that it cannot give the trigger directly, and the drainage is instant, so it has no memory of the rate of flow. The engines of this chapter add the suction from the retention curve and the delay from the conductivity.

\section{Runoff: the curve-number method}
Where rain falls faster than the soil can take it, water runs off. The method of the Soil Conservation Service of the USA estimates the runoff $Q$ from the rain $P$ (mm) of an event by
\begin{equation*}\boxed{Q=\frac{(P-0.2S)^2}{P+0.8S}\quad(P>0.2S),\qquad S=\frac{25\,400}{CN}-254}\end{equation*}
where $S$ is the potential retention (mm) and $CN$ the curve number, from 30 (permeable, covered soil) to 100 (impervious surface). It is empirical and has no time scale; for a CN of 75 and an event of 50 mm, $S=84.7$ mm and $Q=(50-16.9)^2/(50+67.7)=9.3$ mm. The explicit solution of the infiltration (Green--Ampt, Chapter~\ref{chap:11}) is used where the intensity matters.

\section{Stochastic weather: how rain is generated}
When the weather of a future or a different climate is needed, a \kw{weather generator} makes a series of days with the right statistics. Rain is made in two steps. The \emph{occurrence} is a first-order Markov chain: the probability that a day is wet depends on whether the previous day was wet or dry, with two parameters $p_{WW}$ and $p_{DW}$. The \emph{amount} of a wet day follows a skewed distribution, usually the gamma distribution with a shape and a scale parameter, fitted to the record. Temperature and radiation follow with their seasonal means and their dependence on whether the day is wet. The generated series reproduces the mean, the variance and the frequency of wet spells of the record, but not the sequence of the years.

\section{Closure of the budget as a test}
A water-balance model must close: the change of storage equals the inputs less the outputs, to machine accuracy. The \kw{closure residual} is a test of the program, not of the physics. A residual of $10^{-12}$ mm means that no water is created or destroyed in the code; a residual that grows with time points to a bug in one of the fluxes. Section~\ref{sec:closure} shows the closure of the engines.
